\chapter{Extracted Application Kernels}
\label{ch:extracted}
% Data: extracted_kernels_lars/ (copied from HPCAgent-Bench origin/main 299e23df7).

Most benchmarks for loop optimisation are short, synthetic loops. They are good diagnostics,
but they say little about the code that real applications spend their time in. We therefore
extracted eight kernels from six real codes and added them to
HPCAgent-Bench~\cite{hpcagentbench}. This chapter explains which kernels we chose, what they
compute, and how we checked that each port matches its source. \Cref{ch:skills} describes the
extraction procedure itself.

\section{Selection}
\label{sec:extracted-selection}

To organise the choice, we used the thirteen computational patterns, or \emph{dwarfs}, of the
Berkeley view on parallel computing~\cite{asanovic2006landscape}; HPCAgent-Bench groups its
scientific kernels the same way. Our eight kernels cover six of the dwarfs
(\Cref{tab:extracted-kernels}).

For each application, we profiled a realistic run to find where the time goes
(\Cref{tab:profiles}). We used \texttt{perf} for CPU code and Nsight Systems for GPU kernels, on
one GH200 node of CSCS Alps \texttt{daint}. In five of the six applications, one region of
code dominates. QuaTrEx is the exception, which we discuss below.

\begin{table}[t]
  \centering
  \small
  \begin{tabular}{@{}lp{3.5cm}p{3.8cm}r@{}}
    \toprule
    Application & Profiled workload & Hot spot & Share \\
    \midrule
    CoMet          & 2-way CCC, 24\,000 vectors $\times$ 40\,000 fields & CUTLASS INT4 GEMM & 99.7\% GPU \\
    QuaTrEx        & carbon-nanotube GW, 100 energies, 4 iterations (CPU backend) & ZGEMM, incl.\ the RGF block products & $\approx$40\% CPU \\
    SpBench/cuBool & $C = A A$ on five SuiteSparse graphs & row analysis and binning & 55.1\% GPU \\
                   &                                     & hash count and fill     & 44.6\% GPU \\
    WarpX          & 3D plasma, $256^3$ cells, 537\,M particles, 500 steps & field gather (in \texttt{PushPX}) & $\le$42.7\% \\
                   &                                     & current deposition      & 40.8\% \\
                   &                                     & Boris push              & 8.7\% \\
    GraphAIBench   & triangle counting on com-Orkut       & \texttt{triangle\_bs\_warp\_edge} & 100\% GPU \\
    VASim          & Brill, EntityResolution, Fermi automata & per-symbol simulation step & $>$99\% CPU \\
    \bottomrule
  \end{tabular}
  \caption{Hot spots of the upstream applications. GPU shares are fractions of GPU kernel
    time (Nsight Systems), CPU shares fractions of sampled self time (\texttt{perf}).}
  \label{tab:profiles}
\end{table}

\begin{table}[t]
  \centering
  \small
  \begin{tabular}{@{}lll@{}}
    \toprule
    Kernel & Upstream & Data type \\
    \midrule
    \multicolumn{3}{@{}l}{\emph{Dense linear algebra}} \\
    \quad\texttt{comet\_int4\_gemm}  & CoMet~\cite{comet,joubert2019ccc}         & int4 $\to$ int32 \\
    \multicolumn{3}{@{}l}{\emph{Sparse linear algebra}} \\
    \quad\texttt{quatrex\_rgf}       & QuaTrEx~\cite{quatrex,vetsch2025quatrex}     & complex128 \\
    \quad\texttt{spgemm\_hash}       & SpBench/cuBool~\cite{spbench,cubool} & int64 (Boolean) \\
    \multicolumn{3}{@{}l}{\emph{N-body methods}} \\
    \quad\texttt{warpx\_boris\_push} & WarpX~\cite{warpx}         & float64 \\
    \quad\texttt{warpx\_esirkepov\_deposition} & WarpX            & float64 \\
    \quad\texttt{warpx\_field\_gather} & WarpX                    & float64 \\
    \multicolumn{3}{@{}l}{\emph{Graph traversal}} \\
    \quad\texttt{triangle\_count}    & GraphAIBench~\cite{graphaibench,chen2020pangolin} & int64 \\
    \multicolumn{3}{@{}l}{\emph{Finite-state machines}} \\
    \quad\texttt{nfa\_frontier}      & ANMLZoo/VASim~\cite{wadden2016anmlzoo,vasim} & uint8, int64 \\
    \bottomrule
  \end{tabular}
  \caption{The eight extracted kernels, grouped by Berkeley dwarf.}
  \label{tab:extracted-kernels}
\end{table}

The extracted kernels differ from the loop-transformation kernels in three ways. Their memory
accesses are irregular: scatters, hash tables, binary searches and worklists. Most of them do
not compute in \texttt{float64}: only WarpX does, while the others use small integers, bytes or
complex numbers. And several of them branch on the data or on the problem configuration.
Together, the two kernel sets show whether a result holds beyond simple floating-point loops.

\section{Kernel Descriptions}
\label{sec:extracted-descriptions}

Each kernel has a NumPy reference, which defines the correct result, an input generator, a
manifest with its problem sizes, and a native reference that records where the code came from.

\paragraph{CoMet INT4 GEMM.}
CoMet compares genomic vectors in pairs. Its CCC metric is a matrix multiplication over 2-bit
codes, computed on tensor cores: for every pair of vectors $(i, j)$ and every
$(e_i, e_j) \in \{0,1\}^2$,
\[
  \mathit{out}_{i,j,e_i,e_j} = \sum_{f} c_{e_i}(L_{i,f})\, c_{e_j}(R_{j,f}),
\]
where $c_1(v) = \operatorname{popcount}(v)$ and $c_0(v) = 2 - c_1(v)$. We kept this computation
and dropped what only makes sense on a GPU: the tensor-core tiling, the MPI layer, and two small
helper kernels (about 0.1\% of GPU time each). The result is an integer count and must match
exactly.

\paragraph{QuaTrEx RGF selected solve.}
QuaTrEx simulates electron transport through nanoscale devices. Its core is the recursive
Green's function (RGF) algorithm~\cite{svizhenko2002rgf}, which computes selected blocks of the
inverse of a block-tridiagonal matrix without forming the full inverse. This reduces the cost
from $\mathcal{O}((N_B B)^3)$ to $\mathcal{O}(N_B B^3)$ for $N_B$ blocks of size $B$. The
algorithm sweeps forward and then backward over the blocks, once per energy point. The energy
points are independent of each other, but each sweep is strictly sequential. We transcribed it
with six simplifications, all documented in the code, such as using one block size throughout.

QuaTrEx is the only application without a clear hot spot. On the CPU, about 40\% of the run is
complex matrix multiplication, which is what RGF consists of, but the profile cannot tell RGF
apart from the other multiplications in the code. On the GPU, no single kernel takes more than
8.6\% of the time. We chose RGF because it is the algorithmic core of the code, not because the
profile singles it out.

\paragraph{Boolean SpGEMM with hash accumulation.}
This kernel multiplies two sparse Boolean matrices, the operation behind
cuBool~\cite{cubool}, which uses the nsparse algorithm~\cite{nagasaka2017nsparse}. Each row
of the result is the union of some rows of the second matrix, so the work lies in removing
duplicates, not in arithmetic. The original code does this in five phases: it estimates the
size of each row, groups rows of similar size, counts distinct entries with a hash table per
row, computes a prefix sum, and finally fills and sorts the rows. The first two phases take
55\% of the GPU time and the hash phases 45\%, so we ported all five. We left out a special
path for very long rows: it takes 0.1\% of the GPU time and, on one of our test graphs,
produces duplicate entries.

\paragraph{WarpX profile.}
WarpX simulates plasmas with the particle-in-cell method. We profiled a 3D plasma run with
inputs scaled from the WarpX study of Myers et al.~\cite{myers2021warpx}: 537 million particles
on a $256^3$ grid for 500 time steps, which took 16 minutes. Three particle kernels take about
92\% of the time (\Cref{tab:profiles}): gathering the fields to the particles, depositing the
particles' current on the grid, and pushing the particles. The field solver takes only 0.5\%.
The gather is inlined into the push loop, so its 42.7\% also contains part of that loop. We
extracted all three kernels.

\paragraph{WarpX Boris push.}
The Boris pusher~\cite{boris1970} rotates each particle's momentum in the local electric and
magnetic field. Each particle touches only its own data, so the kernel is a simple map, and its
result must match exactly. All three variants of the push are kept.

\paragraph{WarpX Esirkepov current deposition.}
The Esirkepov scheme~\cite{esirkepov2001} writes each particle's current onto the grid such that
charge is conserved exactly. We kept every branch of the original, including all six
geometries. Unlike the other two WarpX kernels, this one is a scatter: neighbouring particles
write to the same grid cells. WarpX uses atomic additions for this, so the order of the
additions, and with it the last bits of the result, changes from run to run. The benchmark
requires identical results across runs, so a parallel version must fix the order of the
additions.

\paragraph{WarpX field gather.}
The field gather does the reverse: it interpolates the fields from the grid to every particle.
The grid is only read, so the particles are independent of each other. The graded run uses a
fixed geometry; \Cref{sec:extracted-design} explains why.

\paragraph{Triangle counting.}
This kernel counts the triangles in a graph, using the binary-search variant of
GraphAIBench~\cite{graphaibench}. Each edge is first directed from the vertex with fewer
neighbours to the one with more, so that every triangle is counted once. For each edge, the
kernel then counts the common neighbours of its two ends, searching the shorter neighbour list
in the longer one. On the com-Orkut graph, it is the application's only GPU kernel and finds
the known 627\,584\,181 triangles in 68\,ms.

\paragraph{NFA frontier simulation.}
This kernel runs a nondeterministic finite automaton over a stream of bytes, as in the ANMLZoo
suite~\cite{wadden2016anmlzoo} and its simulator VASim~\cite{vasim}. In each step, it checks
which active states match the next byte, activates their successors, and re-activates the start
states. An automaton consists of many independent parts, one per pattern, which are the natural
unit of parallel work. This step takes over 99\% of VASim's run time. We left out counter and
gate elements: only one of the twelve ANMLZoo benchmarks uses them, and there most of their cost
is the C++ map they are stored in, not the automaton logic.

\section{Design Decisions and Validation}
\label{sec:extracted-design}

\paragraph{Native references.}
The native reference takes one of three forms. For SpGEMM, triangle counting and the NFA, it
is the original source. For WarpX and CoMet, whose code depends on large libraries, it is a
standalone C++ copy that keeps the original order of operations. For QuaTrEx, which is written
in Python, it is a NumPy copy. The SpGEMM and triangle-counting references are CUDA code, so they
cannot run on a CPU.
% TODO: in the repo, comet_int4_gemm_reference.cpp is translator output and the actual
% extraction is comet_int4_gemm_ref.cpp. Resolve before the evaluation uses it.

\paragraph{Problem sizes.}
Each kernel has four sizes, from S to XL. Where possible, we matched a size to a real input: the
M size of the NFA kernel is close to ANMLZoo's Brill automaton, and the L size of QuaTrEx uses
the block size of a real carbon-nanotube simulation.

\paragraph{Fixed configuration.}
Some kernels also have configuration inputs, such as the geometry of the WarpX kernels. For the
field gather, we fixed the geometry and two other settings at compile time. As run-time inputs,
they kept every branch of the code alive, and the DaCe version exceeded the 1200-second limit of
the benchmark's correctness test. The other geometries are still tested, outside the graded run.

\paragraph{Validation.}
Every port is tested against its native reference, and every generated version against the
NumPy reference. On top of that, we checked each kernel against something that does not depend
on its implementation:
\begin{itemize}
  \item \textbf{QuaTrEx:} the physical symmetries of the result, which a port with a missing
    complex conjugate would break.
  \item \textbf{SpGEMM:} a simple set-union version, and the real cuBool library on
    road-network graphs.
  \item \textbf{Triangle counting:} complete graphs, whose number of triangles is known
    exactly.
  \item \textbf{NFA:} a dense-matrix version of the automaton, and VASim's own statistics on
    four ANMLZoo automata.
  \item \textbf{WarpX and CoMet:} the C++ copies, over all configurations (WarpX) or many
    different inputs (CoMet).
\end{itemize}
Every simplification is documented in the kernel's source, because an undocumented
simplification looks just like a bug.
